\documentclass[10pt,a4paper]{amsart}
\usepackage[margin=19mm]{geometry}
\usepackage{amsmath,amssymb,mathtools}
\usepackage[scr=rsfs,cal=euler]{mathalfa}
\usepackage{xfrac}
\usepackage[unicode,hidelinks,bookmarks=false]{hyperref}
\newcommand{\ie}{i.e.\ }
\newcommand{\cf}{cf.\ }
\newcommand{\resp}{respectively\ }
\newcommand{\Subsection}{\subsection}
\providecommand{\nobreakdash}{\nobreak\hbox{-}}
\setlength{\emergencystretch}{2em}
\multlinegap0pt
\usepackage{amssymb}
\newtheorem{lemma}{Lemma}
\newtheorem{theorem}{Theorem}
\newcommand{\abs}[1]{\left\lvert #1\right\rvert}
\newcommand{\bbc}{\mathbb{C}}
\newcommand{\brac}[1]{\left( #1\right)}
\renewcommand{\geq}{\geqslant}
\renewcommand{\leq}{\leqslant}
\newcommand{\ls}{\lesssim}
\title{More on the sum-product problem for integers with few prime factors}
\author{Rishika Agrawal, Thomas F. Bloom, Giorgis Petridis}
\date{Source: arXiv:2512.04931v2 (CC BY 4.0)}
\begin{document}
\maketitle

\noindent\emph{Source and licence.} More on the sum-product problem for integers with few prime factors. Version arXiv:2512.04931v2 (CC BY 4.0), distributed by arXiv under the Creative Commons Attribution 4.0 International licence, http://creativecommons.org/licenses/by/4.0/, as recorded in the arXiv licence metadata for this exact version and shown on the abstract page https://arxiv.org/abs/2512.04931v2. Full licence notice: \texttt{licenses/arxiv-2512.04931-CC-BY-4.0.txt}.\par\smallskip

\noindent\textit{Editorial scope.} Counted unit: the article's theorem th-growth (source0.tex lines 549--556). The article's complete original proof of it is delivered with it (source0.tex lines 557--576, one \texttt{proof} environment of 20 lines). No mathematics is written, repaired or re-derived: the statement and the proof are the article's own lines, in the article's own notation, and no retained line is substituted or re-flowed.  Retained uncounted as the source-local context the proof needs: lines 362--367; lines 447--450; lines 524--528.  Nothing was omitted from any retained span, so the package carries no omission record.  No retained line contains a citation, so the wrapper carries no bibliography and no bibliography entry.  No retained line contains a figure environment, an image-inclusion command or drawing code, so no image is delivered and the package declares no asset dependency.  Editorial formatting only, in addition: an AMS \texttt{amsart} preamble that restates the article's own declarations and macros used by the retained lines, namely the environments \texttt{lemma}, \texttt{theorem} on separate counters, together with the standard packages those lines need.  The retained environments print in this document as Lemma 1, Theorem 1 in order of appearance, whereas the article prints the same items with the numbering of its own full document.  The complete native source of the article as distributed in the arXiv e-print for this exact version is bundled unchanged as \texttt{sources/190363/source0.tex}.
\begin{lemma}\label{lem-point}
If $r,M\geq 1$ and $\delta\in (0,1)$ are such that 
\[\delta \log M \geq \max(r, (\log M)^{1/6})\]
and $A\subset \bbc^\times$ is $M$-covered by a rank $r$ multiplicative group then, for any $x\neq 0$, 
\[1_A\circ 1_A(x)\ll M^{1+O(\delta^{1/5})}.\]
\end{lemma}

\begin{lemma}\label{lem-energy}
If $A\subset \bbc$ is $M$-covered by a rank $r$ multiplicative group then, for any $m\geq k\geq 1$,
\[E_{2k}(A)\leq 2^{O(km)}\abs{A}^{k}+m^{O(km^3(m+r))}\abs{A}M^{2k-2+\frac{k-1}{m-1}}.\]
\end{lemma}

\begin{lemma}\label{lem-pop}
If either $\abs{A+A}\leq K\abs{A}$ or $\abs{A-A}\leq K\abs{A}$ then, for any $\ell\geq k\geq 1$,
\[\abs{A}^{6k-3}\ll_k K^{2k-1}E_{4\ell}(A)^{\frac{k-1}{\ell-1}}\brac{ K^{2k-1}\abs{A}^{3-2k}E_{4\ell}(A)^{\frac{k-1}{\ell-1}}+ \sum_{x\neq 0}1_A\circ 1_A(x)^{2k}},\]
where the implied constant depends on $k$ only.
\end{lemma}

\begin{theorem}\label{th-growth}
Let $r,M\geq 1$ and $\delta\in (0,1)$ be such that
\[\delta \log M \geq \max(r,(\log M)^{1/6}).\]
If $A\subset \bbc$ is $M$-covered by a rank $r$ multiplicative group and 
\[\min(\abs{A+A},\abs{A-A})=K\abs{A}\]
then
\[\max(K,M)\gg \abs{A}^{5/7-O(\delta^{1/5})} .\]
\end{theorem}

\begin{proof}
By Lemma~\ref{lem-energy}
\[E_4(A) \ll 2^{O(m)}\abs{A}^2+m^{O(m^3(r+m))}\abs{A}M^{2+\frac{2}{m-1}}.\] 
We apply this with $m=\lfloor \delta^{-1/5}\rfloor$, so that
\[E_4(A) \lesssim \abs{A}^2+\abs{A}M^2,\]
where $\ls$ hides losses polynomial in $\abs{A}^{\delta^{1/5}}$. By Lemma~\ref{lem-point} it follows that
\[\sum_{x\neq 0} 1_A\circ 1_A(x)^4\lesssim M^2\brac{\abs{A}^2+\abs{A}M^2}.\]
Applying Lemma~\ref{lem-pop} with $k=2$ therefore implies, for any $\ell\geq 2$,
\[\abs{A}^{9}\ls K^{3}E_{4\ell}(A)^{\frac{1}{\ell-1}}\brac{K^3\abs{A}^{-1}E_{4\ell}(A)^{\frac{1}{\ell-1}}+M^2\abs{A}^2+M^4\abs{A}}.\]
By Lemma~\ref{lem-energy} again 
\[E_{4\ell}(A)^{\frac{1}{\ell-1}} \ls \abs{A}^{2+O(1/\ell)}+\abs{A}^{O(1/\ell)}M^{4}\ls \abs{A}^2+M^4\]
taking $\ell$ sufficiently large ($\asymp \log \abs{A}$ say). It follows that 
\[\abs{A}^{9}\ls K^{3}(\abs{A}^2+M^4)\brac{K^3\abs{A}^{-1}(\abs{A}^2+M^4)+M^2\abs{A}^2+M^4\abs{A}}.\]
Simplifying the right-hand side yields
\[\abs{A}^9\ls K^6\abs{A}^3+K^6M^4\abs{A}+K^3M^2\abs{A}^4+K^3M^4\abs{A}^3\]
\[ +K^6M^8\abs{A}^{-1}+K^3M^6\abs{A}^2+K^3M^8\abs{A}\]
which implies
\[\max(K,M) \gtrsim \abs{A}^{5/7}\]
as claimed. 
\end{proof}

\end{document}
